\subsection{The unit eighth moments for minimum-three even lists and the pair $(2,\allowbreak 3)$}\label{proof:Q5-C053-unit}
Use $A=\allowbreak J+K$, $T=\allowbreak J-K$, $g=\allowbreak K\one$, $p=\allowbreak 1+g$, $q=\allowbreak 1-g$,
$d=\allowbreak \E g^2$, $\delta_j=\allowbreak \|K^j\one\|^2$, $\theta=\allowbreak 13/32$, $\kappa=\allowbreak 19/32$,
$f_n=\allowbreak A^n\one$ and $h_n=\allowbreak Tf_{n-1}$.
For $3=\allowbreak b_1<\allowbreak c_1<\allowbreak b_2<\allowbreak c_2<\allowbreak \cdots<\allowbreak b_j<\allowbreak c_j$, set
\[
 F=f_3h_{c_1}\prod_{i=2}^j h_{b_i}f_{c_i}.
\]
We prove
\begin{equation}\tag{Q53.1.4}\label{eq:Q53:minthree-unit}
 \E(pTF)^8\le e^{-(16j-61/4)d}\le e^{-3d/4}.
\end{equation}
We also prove
\begin{equation}\tag{Q53.1.3}\label{eq:Q53:twothree-unit}
 \E[pT(h_2f_3)]^8\le1-d/500.
\end{equation}
In these formulas the eighth power is inside the expectation. All directions
are determined by the listed arm lengths. The receiving unit arm points
outward, and the unit connector points from the terminal toward that root.

\paragraph{Q53.2 Exponential transfer at a unit connector}\label{part:Q53.2}
For $|u|\le\allowbreak 8/5$,
\begin{equation}\tag{Q53.2.1}\label{eq:Q53:small-exp}
 e^u\le1+u+u^2.
\end{equation}
Divide the absolute Taylor tail by $u^2$. Its quadratic term is $1/2$,
the cubic term is at most $4/15$, and the fourth and later terms sum to at
most $(8/75)/(1-8/25)=\allowbreak 8/51$. Their total is $157/170<\allowbreak 1$.
Suppose $F\ge\allowbreak 0$, $F\le\allowbreak  e^{R-\Lambda d}$, $|R|\le\allowbreak 8/5$ and $|8KR|\le\allowbreak 8/5$.
Write $V=\allowbreak \E R^2$ and $Z=\allowbreak KR$. Since the kernel of $T$ is at most two,
\[
 Te^R\le q+\E R-Z+2V=:U.
\]
The function $U$ is nonnegative. Arithmetic--geometric mean gives
\[
 pTF\le e^{-\Lambda d}pU
 \le e^{-\Lambda d}[1+\E R/2+V-Z/2]^2.
\]
The base is $(p+U)/2\ge\allowbreak 0$. Raising to the eighth power and using
$1+y\le\allowbreak  e^y$, followed by \eqref{eq:Q53:small-exp} for $-8Z$, yields
\begin{equation}\tag{Q53.2.2}\label{eq:Q53:transfer}
 \E(pTF)^8\le
 e^{-8\Lambda d+8\E R+16V-8\E Z+64\E Z^2}.
\end{equation}
Because $\E Z^2\le\allowbreak \theta V$ and $16+64\theta=\allowbreak 42$, the conditions
$\E R\le\allowbreak  md$, $V\le\allowbreak  vd$, $-\E Z\le\allowbreak  ad$ imply
\begin{equation}\tag{Q53.2.4}\label{eq:Q53:budget}
 \E(pTF)^8\le e^{-(8\Lambda-8m-42v-8a)d}.
\end{equation}
If instead $F\le\allowbreak  e^{-cd}$ for $c\ge\allowbreak 0$, positivity and $pq\le\allowbreak 1$ give
\begin{equation}\tag{Q53.2.5}\label{eq:Q53:pointwise}
 \E(pTF)^8\le e^{-8cd}.
\end{equation}

\paragraph{Q53.5 The minimum-three list}\label{part:Q53.5}
The retained envelope of Section~\ref{loc:S28}, with the same direction
$f_3h_b$, gives a function $R_0$ such that
\begin{equation}\tag{Q53.5.1}\label{eq:Q53:firstpair}
 f_3h_b\le e^{R_0-3d},\qquad \E R_0\le-\kappa\delta_2.
\end{equation}
Its explicit estimates \eqref{eq:compact-min3-small} and
\eqref{eq:compact-min3-large} are
\begin{equation}\tag{Q53.5.2}\label{eq:Q53:pair-ranges}
 \begin{cases}
 \|R_0\|\le(31/100)\sqrt d,\quad |R_0|\le1/4,&0<d\le1/9,\\
 |R_0|\le(69/32-1/4000)d,&1/9\le d\le1/3.
 \end{cases}
\end{equation}
For the mean, the spectral representation there is
$\E R_0=\allowbreak -\int B(x)\,\allowbreak d\nu_g$, where $B=\allowbreak D+xb_x$, $D\ge\allowbreak 0$ and
$b_x\ge\allowbreak 1-x\ge\allowbreak \kappa$.
Set $S=\allowbreak \sum_{i=\allowbreak 2}^jR_{b_ic_i}$, with $S=\allowbreak 0$ for $j=\allowbreak 1$.
Each later shorter arm is at least five, and each of its exponents begins
at six. The finite-mask bound \eqref{eq:L90:mask} and the pair estimate give
\begin{equation}\tag{Q53.5.3}\label{eq:Q53:whole-envelope}
 F\le e^{R-(2j+1)d},\qquad R=R_0+S.
\end{equation}
Writing $S=\allowbreak K^5P(K)g$ with a finite zero--one polynomial gives
\begin{equation}\tag{Q53.5.4}\label{eq:Q53:tail}
 \|S\|\le(3/2)\theta^{5/2}\sqrt d<\sqrt d/6,
 \qquad |S|\le(3/2)\theta^2\sqrt d<\sqrt d/4.
\end{equation}
Here $\theta^2<\allowbreak 1/6$ and $\sqrt\theta<\allowbreak 2/3$. Its possible positive mean
starts at $\delta_4$, so
\begin{equation}\tag{Q53.5.5}\label{eq:Q53:mean}
 \E S\le\frac{\delta_4}{1-\theta^2}
 \le\frac{\theta^2}{1-\theta^2}\delta_2<\delta_2/5,
 \qquad\E R\le-(\kappa-1/5)\delta_2\le0.
\end{equation}
For $0<\allowbreak d\le\allowbreak 1/9$,
\[
\begin{gathered}
\|R\|\le(143/300)\sqrt d<\sqrt d/2,\quad\\
V\le d/4,\quad |R|\le1/4+\sqrt d/4\le1/3,\\
\quad |8KR|\le4\sqrt d\le4/3<8/5.
\end{gathered}
\]
Skew-adjointness gives $-\E(KR)=\allowbreak \E(gR)\le\allowbreak  d/2$.
Taking $\Lambda=\allowbreak 2j+1$, $m=\allowbreak 0$, $v=\allowbreak 1/4$ and $a=\allowbreak 1/2$ in
\eqref{eq:Q53:budget} gives
\begin{equation}\tag{Q53.5.6}\label{eq:Q53:small-payment}
 \E(pTF)^8\le e^{-(16j-13/2)d}.
\end{equation}
For $1/9\le\allowbreak  d\le\allowbreak 1/3$, use $\sqrt d\le\allowbreak 3d$ to obtain
$R\le\allowbreak (93/32)d$, and hence
\begin{equation}\tag{Q53.5.7}\label{eq:Q53:large-envelope}
 F\le e^{-(2j-61/32)d}.
\end{equation}
The loss is at least $3/32$ when $j\ge\allowbreak 1$, so
\eqref{eq:Q53:pointwise} gives
\begin{equation}\tag{Q53.5.8}\label{eq:Q53:large-payment}
 \E(pTF)^8\le e^{-(16j-61/4)d}.
\end{equation}
The small-domain bound is stronger, proving \eqref{eq:Q53:minthree-unit}.

\paragraph{Q53.7 The pair $(2,\allowbreak 3)$}\label{part:Q53.7}
Put $r=\allowbreak Kg$, $s=\allowbreak Kr$, $t=\allowbreak Ks$ and $a=\allowbreak 1-d/2\ge\allowbreak 5/6$.
The identities $h_2=\allowbreak q-r$, $f_3=\allowbreak p-d+r+s$ give
$h_2+f_3=\allowbreak 2-d+s$, so $h_2f_3\le\allowbreak (a+s/2)^2$.
Since $\E s=\allowbreak 0$, $Ts=\allowbreak -t$ and $T(s^2)\le\allowbreak 2\delta_3$,
\[
 T(h_2f_3)\le a^2q-at+\delta_3/2=:U\ge0.
\]
Arithmetic--geometric mean for $ap$ and $U/a$ gives
\begin{equation}\tag{Q53.7.1}\label{eq:Q53:twothree-square}
 0\le pT(h_2f_3)\le[a-t/2+\delta_3/(4a)]^2.
\end{equation}
The base is half the sum of the two nonnegative factors. Moreover,
\[
 a+\delta_3/(4a)\le1-d/2+(3/10)\theta^2d
 =1-(4613/10240)d\le1-(51/128)d=:c.
\]
The two loss coefficients differ by $533/10240>\allowbreak 0$. Thus
\begin{equation}\tag{Q53.7.3}\label{eq:Q53:twothree-majorant}
 0\le pT(h_2f_3)\le(c-t/2)^2.
\end{equation}
The row estimate of Section~\ref{proof:Q5-C049}, applied to $t=\allowbreak K^2r$, gives
\begin{equation}\tag{Q53.7.4}\label{eq:Q53:twothree-data}
 \E t=\delta_2\ge0,\qquad |t|^2\le\delta_2/3\le d/3,
 \qquad \E t^2=\delta_4\le\theta^3d\le\theta^2d.
\end{equation}
These are precisely the hypotheses of \eqref{eq:Q51:2.2}; with $Z=\allowbreak t$ it
proves \eqref{eq:Q53:twothree-unit}. Positivity restores all equal-arm
decorations. At $d=\allowbreak 0$ the bounds follow by leaf integration, and replacing
$K$ by $-K$ gives their full reversals.
